\documentclass[output=paper,colorlinks,citecolor=brown,
% hidelinks,
% showindex
]{langscibook}

%This is where you put the authors and their affiliations
\author{We'll insert your name here\affiliation{We'll insert your affiliation here}}

%Insert your title here
\title{ACAL Proceedings LaTeX-free Template: We'll insert your title here for you}  
\abstract{Here is where we will insert the abstract of the paper for you}

\input{AcalProceedingsTemplatePackages.tex}
\input{AcalProceedingsTemplateCommands.tex}

% \IfFileExists{../localcommands.tex}{%hack to check whether this is being compiled as part of a collection or standalone
%   \input{../localpackages}
%   \input{../localcommands} 
% \togglepaper[23]
% }{}

\begin{document}
\maketitle


\section{What is this document?}

BEFORE YOU DO ANYTHING ELSE: In the header on this page there is a Green Button that says ``Source,'' paired with an alternative that says ``Rich Text.'' Please click on ``Rich Text'' to toggle your view from Source to Rich Text. 

OK, now proceed. :-) 

This is a template for submitting for submitting proceedings papers in a LaTeX-compatible format for the Annual Conference on African Linguistics, published as an Overleaf template. It does not require LaTeX knowledge, but is designed to make later conversions much easier.

The purpose of this template is for people who are using the ACAL proceedings hybrid format of submission. The purpose of this is to make it easier for editors to convert your paper to LaTeX, by you doing some of the basic work and leaving the more complex formatting to them. This submission strategy consists of a few main components: 
\begin{enumerate}
\item You will copy the main text of your document into this Overleaf project. 
\item You will put your bibliography into a LaTeX-compatible format
\item You will submit examples/trees/charts/images/etc separately in BOTH PDF format and whatever format you initially created them in. 
\item Be sure to submit a the Word and PDF versions of your original document before you did the conversion here. This will help us know what your desired outcome looks like. 

\end{enumerate}
Full instructions for how to do each stage of this are posted on the ACAL website under ``Annual Conference Proceedings.'' Here, we just focus on submitting the main text of your paper into this document.

\section{How to use this template}
\begin{itemize}

\item If you opened this document on LaTeX and signed into Overleaf, a copy of this project is now in your Overleaf account (i.e. it belongs to you), and you can now edit this document to be an actual paper that you write. 
\item This means there is no setup required, you can just erase this text and start typing your own paper, or use this document to practice and see what changes your edits make.
\item This also means you will probably want to rename the document - you can do this by hovering your cursor over the title `ACAL proceedings "LaTeX-free" submission template' at the top of this page, and then clicking on the pencil icon that appears to the right of the title. You can also rename the AcalProceedingsRtfTemplate.tex file in the leftmost panel by clicking on the down arrow.
\end{itemize}

\section{Using Overleaf in Rich Text}
\begin{itemize}

\item In Overleaf in Rich Text formatting, you edit the text in the frame that is in the middle of your screen. You will note as you do this, it doesn't automatically change the  document to the right.
\item Pressing the Green ``Recompile" button above the right-side document then converts what you wrote into Language Science Press's formatting. 
\item We recommend doing this relatively frequently while writing, just to make sure mistakes/errors didn't arise. If you do run into errors that you don't know how to fix, you can contact us at acal.latex@gmail.com. 
\item The instructions document on the ACAL website (which you likely followed to get here) has some details about what specific buttons on this page do, e.g. bold, italics, sections, subsections, bulleted lists, etc.
\item Footnotes are written directly into the text at the point where the footnote appears. They are introduced by backslash-footnote, and then the footnote text surrounded by curly brackets, as displayed here.\footnote{Hey look at me, I'm a footnote.} When you Press Recompile, you will see the footnote text appear as a footnote.
\item Please do not insert a references list. Instead, you will compile a list of sources that you cite in a LaTeX-compatible format, as outlined in the document posted on the ACAL website.
\end{itemize}

\section{A couple things to keep an eye out for}

\begin{itemize}
\item Quotation marks work slightly different in this application. So you can see that if I copy them in the way I would have in "Microsoft Word" that the front quotes face the wrong way in the typeset document to the right. Instead, front quotes use the grave accent mark ` which appears at the left of your keyboard above the tab key. ``like this''
\item If you need to use one of the following symbols:  \$, \&, \%  you will need to put a backslash in front of it, as we did here. The reason is that inside LaTeX those symbols mean something specific, so if you want them to just mean what they mean in normal life (i.e. dollar symbol, `and', percentage), the backslash essentially says ``hey I just want an actual percentage symbol.”  
\end{itemize}


\section{Want to ``level up'' your submission to be extra helpful to the editors?}

There are several small LaTeX formatting skills you can pick up which would assist the volume editors a lot. None are required to submit to ACAL, but all are relatively easy to learn if you wanted to. Here we look at formatting data examples, formatting cross-references, and inserting citations.

But, to be clear: if even looking at anything below makes your head spin, you are welcome to ignore this. 

\subsection{You can learn to format data examples}

First, you can learn to format data examples. LaTeX makes this relatively easy - it auto-aligns glosses, for example. One you learn the basic framework for inserting data examples LaTeX will do all the work for you. 

First, a numbered example alone is inserted using the code backslash-ea to start the example, and backslash-z to end the example.

\ea 
This is a numbered example.
\z

This also allows for inter-linear glossing and translations: 

\ea 
\gll this is a language example \\
this is the morphological gloss \\
\glt `This is the translation.' 
\z 

This will automatically aligns glosses with the language example (each white space means a different word), so if you press the Green Recompile Button and scroll to this position in the paper, you will see the aligned gloss . This is useful for non-English examples like the Lubukusu sentence in the next example

\ea 
\gll Wekesa se-a-la-ba a-kula ka-ma-indi ta. \\
Wekesa NEG-SM-PST-be SM-buy 6-6-maize NEG \\ 
\glt `Wekesa will not be buying maize.' (Justine Sikuku, pc)
\z

Some things to notice for the examples above - an example is started with the backslash-ea and ended with backslash-z - omitting either of these will generate an error. The language and gloss lines are started with the annotation backslash-gll, with line breaks (double backslashes) ending both the language line and the gloss line. The translation line is introduced with the annotation backslash-glt, and notice that there is no line break at the end of the translation line.

If you want multiple data examples to appear in one numbered example, you use the xlist environment to create a sublist. 

\ea
\begin{xlist}

\ex This is the first example. 

\ex 
\gll hii ni m-fano w-a pili, kwa Ki-swahili \\
this is 3-example 3-ASSOC second, of 7-Swahili \\
\glt `This is a second example, in Swahili.' 

\end{xlist}

\z 

\noindent Again, some details to notice: the annotation backslash-begin{xlist} and backslash-end{xlist} that begin and end the sublist inside the example. It is important that the sublist be generated \textit{inside} the numbered example, so backslash-begin{xlist} comes after backslash-ea and backslash-end{xlist} comes before backslash-z. Each example inside the sublist is introduced with the annotation backslash-ex. 

\subsection{You can learn to format cross-references}

If you want to format cross-references, there are two components to keep in mind. First, you create a label on the example you want to refer to, and then you create the cross-reference. A label is created with the annotation backslash-label{label.name}. Language Science Press asks us to choose label names that are partially predetermined, and partially chosen by us. So in the example below, \verb|\label{ex:diercks:XlistExample}| is the label of the entire example, and \verb|\label{ex:diercks:SwahiliExample}| is the label of the Swahili example.  As above, please ignore the \verb|| portion of the annotation above, that is just to allow you to see the whole code.  As you will see,  retype (or copy-paste) just the red portion of the preceding annotations inside the document, and Overleaf will convert it to the ``label'' icon you see below. 

\ea \label{ex:diercks:XlistExample}
\begin{xlist}

\ex This is the first example. 

\ex \label{ex:diercks:SwahiliExample}
\gll hii ni m-fano w-a pili, kwa Ki-swahili \\
this is 3-example 3-\Assoc{} second, of 7-Swahili \\
\glt `This is a second example, in Swahili.' 

\end{xlist}

\z 

To make a cross-reference, then, we simply refer to the label we gave to an example. We use the annotation backslash-REF{label.name} (where capital REF is important) to create the cross-reference to the example. So I can say that inside example \REF{ex:diercks:XlistExample}, there is a Swahili example in \REF{ex:diercks:SwahiliExample}. And as you can see in the typeset document on the right

Be sure to choose a different label name for every example, so LaTeX doesn't get confused.

\subsection{You can learn to format in-text citations}

Doing this first requires building your bibliography in a LaTeX-compatible format: those instructions are explained in our documentation on building your bibliography (LINKED). 

After doing that, adding citations to your text is relatively straightforward.

The two most important citation annotations are backslash-citet{} and backslash-citep{}.  The cite key of the reference goes between the curly brackets. (What a `cite key' is is explained in the instructions (LINKED) mentioned above.) You can ignore the \verb|| portion of the code below and focus on the portion that appears in red.
 
 \begin{itemize}
 \item For `Author (Date)' format, use backslash-citet{}. So writing \verb|\citet{Lahiri2000}| produces \citet{Lahiri2000}, which Overleaf's Rich Text reformats to look like a book icon+citekey. 
 \item For a parenthetical citation, use backslash-citep{}. So writing \verb|\citep{Lahiri2000}| produces \citep{Lahiri2000}, which Overleaf's Rich Text also reformats to look like a book icon+citekey. 
 \end{itemize}

\noindent Both of these commands can be modified to add extra information. For some reason, Overleaf's Rich Text \textit{doesn't} add the fancy book icon to these references (which means we don't have to add the \verb|| annotation to show you what to type).
 
 \begin{itemize}
 \item To add a page number after the year, put the page number in square brackets \textit{before} your curly brackets, e.g. \citep[87]{Doke1923} 
 \item To add information before the citation, put this in square brackets before the brackets with the page number, e.g. \citep[Zulu,][87]{Doke1923} 
  \item If you want information before the citation, but without a page number, just leave the page number blank, e.g. \citep[Zulu,][]{Doke1923}
  \item Multiple cite keys separated by commas will produce multiple citations, e.g. \citep{Doke1923,Lahiri2000}
 \end{itemize}
 
This \href{http://merkel.texture.rocks/Latex/natbib.php}{helpful reference sheet} gives many details about citation commands for additional variants that you may want.


\section{Conclusion}


\section*{Abbreviations}

Abbreviations should be included in this section. It is a required part of volumes with Language Science Press that an abbreviations list be included, so you cannot omit this section.

\section*{Acknowledgements}

Acknowledgements must be listed here, not in a footnote in the paper.

\section{Please ignore everything below this note!}

\sloppy %this command eliminates a quirk that appears in the bibliography, you can just leave it here.
\printbibliography[heading=subbibliography,notkeyword=this]

\end{document}
